\section{El patrón namespace para sistemas multi-robot}

\subsection{El problema del nombre único}

Cuando dos robots ejecutan el mismo nodo de control, ambos intentan publicar en /cmd\_vel y suscribirse a /odom. DDS conecta a todos los publicadores y suscriptores del mismo tópico entre sí, sin distinción. Robot 0 recibiría la odometría de robot 1 y viceversa; los comandos de velocidad se mezclarían entre ambos. El sistema colapsaría.

La solución es el namespace: un prefijo jerárquico que califica todos los nombres de tópicos de un nodo. Un nodo con namespace robot0 que publica en cmd\_vel en realidad publica en /robot0/cmd\_vel. Un nodo con namespace robot1 publica en /robot1/cmd\_vel. Los dos canales son distintos y no interfieren.

\subsection{Asignación de namespace en el launch file}

El namespace se asigna mediante el parámetro namespace de la acción Node. Todos los tópicos, servicios, parámetros y el propio nombre del nodo quedan bajo ese prefijo automáticamente, sin necesidad de modificar el código del nodo. El Código~\ref{lst:ns_node} muestra el patrón:

\begin{figure*}[t]
\centering
\begin{lstlisting}[style=python,
    caption={Asignación de namespace a un nodo desde el launch file.},
    label={lst:ns_node}]
for i in range(n):
    ns = f'robot{i}'   # 'robot0', 'robot1', ...

    Node(
        package='control_nodes',
        executable='consenso_node',
        name='consenso',
        namespace=ns,           # <- aqui
        parameters=[{...}],
        output='screen'
    )
\end{lstlisting}
\end{figure*}


Con namespace='robot0' y name='consenso', el nombre completo del nodo en el grafo ROS~2 es /robot0/consenso. Sus tópicos quedan en /robot0/cmd\_vel, /robot0/odom, etc.

\subsection{Namespace en IncludeLaunchDescription}

El mismo principio aplica cuando se incluye el launch file de descripción del robot. El namespace debe pasarse como argumento al launch hijo para que el RSP publique el tópico robot\_description bajo el namespace correcto:

\begin{figure*}[t]
\centering
\begin{lstlisting}[style=python,
    caption={Namespace propagado al launch de descripción del robot.},
    label={lst:ns_rsp}]
rsp_node = IncludeLaunchDescription(
    PythonLaunchDescriptionSource(rsp_launch),
    launch_arguments={
        'namespace':    ns,
        'use_sim_time': 'false'
    }.items()
)
\end{lstlisting}
\end{figure*}


Internamente, rsp.launch.py usa ese argumento para asignar el namespace al nodo robot\_state\_publisher, que entonces publica en /robot0/robot\_description. El nodo de spawn lee ese tópico exacto mediante el flag -topic /{ns}/robot\_description, cerrando la cadena.

\subsection{Verificación desde la terminal}

Con un sistema multi-robot corriendo, la CLI de ROS~2 permite inspeccionar la separación por namespaces desde la Terminal~\ref{lst:ns_inspect}:

{\renewcommand{\lstlistingname}{Terminal}
\begin{figure*}[t]
\centering
\begin{lstlisting}[style=bash,
    caption={Inspección de nodos y tópicos con namespaces.},
    label={lst:ns_inspect}]
# Lista todos los nodos activos
ros2 node list
# /robot0/consenso
# /robot1/consenso
# ...

# Lista topicos de un robot especifico
ros2 topic list | grep robot0
# /robot0/cmd_vel
# /robot0/odom
# /robot0/scan
\end{lstlisting}
\end{figure*}
}

